\subsubsection{Higher-Order Equations and Systems}\label{subsec:higher-order-systems}

The generalization of the definitions above to equations of higher order than second is straightforward.

\begin{Definition}
\label{def:higher-order-pde-type}
Let $L$ be the $m$th-order operator defined in \eqref{eq:general-linear-differential-expression}. Characteristic surfaces are defined by the equation
\begin{equation}\label{eq:higher-order-characteristic-surface}
L^p(\bm{x},\nabla\phi)=0.
\end{equation}
An equation is called \textbf{elliptic} at $\bm{x}$ if there are no real characteristics at $\bm{x}$ or, equivalently, if
\begin{equation}\label{eq:higher-order-ellipticity}
L^p(\bm{x},i\bm{\xi})\ne0,\qquad\forall\bm{\xi}\ne0.
\end{equation}
An equation is called \textbf{strictly hyperbolic}\footnote{We shall not give a general definition of what it means to be nonstrictly hyperbolic, although such definitions exist. Below we shall define nonstrict hyperbolicity for first-order systems.} in the direction $\bm{n}$ if
\begin{enumerate}
\item $L^p(\bm{x},i\bm{n})\ne0$, and
\item all the roots $\omega$ of the equation
\begin{equation}\label{eq:strict-hyperbolic-characteristic-roots}
L^p(\bm{x},i\bm{\xi}+i\omega\bm{n})=0
\end{equation}
are real and distinct for every $\bm{\xi}\in\mathbb{R}^n$ which is not collinear with $\bm{n}$.\end{enumerate}
\end{Definition}

In applications, $\bm{n}$ is usually a coordinate direction associated with time. In this case, let us set $\bm{x}=(x_1,x_2,\ldots,x_{n-1},t)$ and let $\bm{\xi}=(\xi_1,\ldots,\xi_{n-1},0)$ be a spatial vector.

For rapidly oscillating functions of small support, we may think of the coefficients of $L^p$ as approximately constant; let us assume they are constant. If $\omega$ is a root of \eqref{eq:strict-hyperbolic-characteristic-roots}, then $u=\exp(i(\bm{\xi}\cdot\bm{x})+i\omega t)$ is a solution of $L^pu=0$. If $\omega$ has negative imaginary part, then this solution grows exponentially in time. Moreover, since $L^p$ is homogeneous of degree $m$, i.e., $L^p(\bm{x},\lambda(i\bm{\xi}+i\omega\bm{n}))=\lambda^mL^p(\bm{x},i\bm{\xi}+i\omega\bm{n})$ for any scalar $\lambda$, there are always roots with negative imaginary parts if there are any roots which are not real (if we change the sign of $\bm{\xi}$, we also change the sign of $\omega$). Moreover, if we multiply $\bm{\xi}$ by a scalar factor $\lambda$, then $\omega$ is multiplied by the same factor, and hence solutions would grow more and more rapidly the faster they oscillate in space. The condition that the roots in \eqref{eq:strict-hyperbolic-characteristic-roots} are real is therefore a necessary condition for well-posedness of initial-value problems.

We now turn our attention to systems of $k$ partial differential equations involving $k$ unknowns $u_j$, $j=1,2,\ldots,k$:
\begin{equation}\label{eq:linear-pde-system}
L_{ij}(\bm{x},D)u_j=0,\qquad i=1,2,\ldots,k.
\end{equation}
As for systems of algebraic equations, well-posed problems require equal numbers of equations and unknowns, so we shall assume that the operators $L_{ij}$ form a square matrix $\bm{L}$. The generalization of the notions above is in principle quite straightforward.

\begin{Definition}
\label{def:system-pde-type}
Characteristic surfaces are defined by the equation
\begin{equation}\label{eq:system-characteristic-surface}
\det\bm{L}^p(\bm{x},\nabla\phi)=0,
\end{equation}
and equations without real characteristic surfaces are called \textbf{elliptic}. \textbf{Strict hyperbolicity} is also defined as above, with $L^p$ in the definition replaced by $\det\bm{L}^p$. A system in which all components of $\bm{L}^p$ are operators of first order is called hyperbolic (not necessarily strictly) in the direction $\bm{n}$ if
\begin{enumerate}
\item $\det\bm{L}^p(\bm{x},\bm{n})\ne0$, and
\item for $\bm{\xi}$ not collinear with $\bm{n}$, all eigenvalues $\omega$ of the problem
\[
\det\bm{L}^p(\bm{x},i\bm{\xi}+i\omega\bm{n})=0
\]
are real and there is a complete set of eigenvectors.\end{enumerate}
\end{Definition}

Note that, since we assumed that the components of $\bm{L}^p$ are of first order, we have $\bm{L}^p(\bm{x},i\bm{\xi}+i\omega\bm{n})=\bm{L}^p(\bm{x},i\bm{\xi})+\omega\bm{L}^p(\bm{x},i\bm{n})$; hence the problem $\det\bm{L}^p=0$ is a matrix eigenvalue problem for $\omega$. If the eigenvalues are distinct, there is always a complete set of eigenvectors; hence strict hyperbolicity implies hyperbolicity.

In general, we need to be careful about defining the "principal part" of a system. A naive approach of simply taking the "terms of highest order" turns out to be unsatisfactory.

\begin{Example}
\label{ex:weighted-laplace-reformulation}
To see the problem, let us consider Laplace's equation in two dimensions
\begin{equation}\label{eq:first-order-laplace-reformulation-origin}
u_{xx}+u_{yy}=0,
\end{equation}
and rewrite it as a system of first-order equations by setting $v=u_x$, $w=u_y$. The resulting system is
\begin{equation}\label{eq:first-order-laplace-system}
u_x=v,\qquad u_y=w,\qquad v_x+w_y=0.
\end{equation}
If we define $\bm{L}^p$ to be the part involving first-order terms, it is easy to see that $\det\bm{L}^p$ turns out to be identically zero. On the other hand, since Laplace's equation is the standard example of an elliptic equation, it would be desirable to have the equivalent first-order system also defined as "elliptic." Obviously, we then cannot throw away the terms $v$ and $w$ in the first two equations of \eqref{eq:first-order-laplace-system}.
\end{Example}

The difficulty is resolved by assigning "weights" $s_i$ to each equation and $t_j$ to each dependent variable in such a way that the order of each operator $L_{ij}$ does not exceed $s_i+t_j$. The principal part $L_{ij}^p$ is then defined to consist of those terms which have order exactly equal to $s_i+t_j$. We assume that the weights can be assigned in such a way that $\det\bm{L}^p$ does not vanish identically\footnote{There are examples where this assumption fails, e.g., the system $u_x+v_y=0$, $u_x+v_y+v=0$. The difficulty there disappears if we use the equivalent form $u_x+v_y=0$, $v=0$. It is known that weights with the desired properties always exist for nondegenerate systems \cite{Vo}. Here nondegenerate means the following: If $\det\bm{L}(i\bm{\xi})$ is expressed in the usual way as a sum of products, then the degree of each of these products as a polynomial in $\bm{\xi}$ does not exceed the degree of the determinant.}; in this case $\det\bm{L}^p$ consists of all the terms of order $\sum_{i,j}s_i+t_j$ which appear in $\det\bm{L}$. In \eqref{eq:first-order-laplace-system}, for example, we would set $s_1=s_2=t_2=t_3=0$ and $t_1=s_3=1$. (Here, it is understood that the ordering of the variables is $u,v,w$.) With these weights, the principal part of \eqref{eq:first-order-laplace-system} is actually identical to \eqref{eq:first-order-laplace-system}, and we compute
\begin{equation}\label{eq:weighted-laplace-principal-determinant}
\det\bm{L}^p(i\bm{\xi})=\det\begin{pmatrix}i\xi_1&-1&0\\i\xi_2&0&-1\\0&i\xi_1&i\xi_2\end{pmatrix}=-\xi_1^2-\xi_2^2,
\end{equation}
which is equal to the symbol of Laplace's equation.

The term "order" for systems is used with two different meanings. Such terms as "first-order" or "second-order" systems usually refer to the equations of which the system is composed. However, it is also possible to assign an order to the system as a whole. This order is defined as the degree of the symbol of $\det\bm{L}^p$, which is equal to $\sum_{i,j}s_i+t_j$.

\begin{Remark}
\label{rem:nonunique-principal-weights}
We note that weights need not be unique. The following is a simple example of a system which is elliptic under more than one choice of weights. Both choices are useful, and lead to different elliptic regularity results. Consider the system
\begin{equation}\label{eq:two-weight-elliptic-system}
\Delta u-v=0,\quad\Delta v=0.
\end{equation}
We can choose the weights $s_1=s_2=2$, $t_1=t_2=0$. With these choices, the principal part of the symbol is
\begin{equation}\label{eq:first-choice-weighted-principal-symbol}
\bm{L}^p(i\bm{\xi})=\begin{pmatrix}-|\bm{\xi}|^2&0\\0&-|\bm{\xi}|^2\end{pmatrix}.
\end{equation}
But we can also choose the weights $t_1=2$, $t_2=0$, $s_1=0$, $s_2=2$. Now the principal part is
\begin{equation}\label{eq:second-choice-weighted-principal-symbol}
\bm{L}^p(i\bm{\xi})=\begin{pmatrix}-|\bm{\xi}|^2&-1\\0&-|\bm{\xi}|^2\end{pmatrix}.
\end{equation}
In both cases, we have $\det\bm{L}^p(i\bm{\xi})=|\bm{\xi}|^4$.
\end{Remark}

\subsubsection{Nonlinear Equations}\label{subsec:nonlinear-classification}

For nonlinear equations and systems, the type can depend not only on the point in space but on the solution itself. We simply linearize the equation at a given solution and define the type to be that of the linearized equation. Characteristic surfaces are similarly defined as the characteristic surfaces of the linearized equation.

For future use we give the definition of quasilinear and semilinear:

\begin{Definition}
\label{def:quasilinear-semilinear}
A system is called \textbf{quasilinear} if derivatives of principal order occur only linearly (with coefficients which may depend on derivatives of lower order). It is called \textbf{semilinear} if it is quasilinear and the coefficients of the terms of principal order depend only on $\bm{x}$, but not on the solution.
\end{Definition}

\begin{Example}
\label{ex:nonlinear-classification}
The equation
\begin{equation}\label{eq:quasilinear-type-example}
\alpha(u_x)u_{xx}+u_{yy}=0
\end{equation}
is quasilinear; it is elliptic if $\alpha(u_x)>0$ and hyperbolic if $\alpha(u_x)<0$. The equation
\begin{equation}\label{eq:nonquasilinear-type-example}
\phi(u_{xx})+u_{yy}=0
\end{equation}
is not quasilinear; it is elliptic if $\phi'(u_{xx})>0$ and hyperbolic if $\phi'(u_{xx})<0$. The equation
\begin{equation}\label{eq:semilinear-elliptic-example}
(x^2+1)u_{xx}+u_{yy}+\phi(u_x,u_y)=0
\end{equation}
is semilinear elliptic.
\end{Example}

\paragraph*{Problems}

\begin{enumerate}[label=\thesection.\arabic*.,ref=\thesection.\arabic*,start=1,series=book-problems]

\item\label{ex:determine-pde-types} Determine the type of the following equations:\begin{enumerate}[label=(\alph*)]
\item $u_{xy}+2u_x=0$.
\item $u_{xxxx}-u_{xxyy}+u_{yyyy}=0$.
\item $u_{tt}+u_{xxxx}=0$.
\end{enumerate}

\item\label{ex:find-second-order-characteristics} Find the characteristics of the following equations.\begin{enumerate}[label=(\alph*)]
\item $xu_{xx}+u_{yy}=0$.
\item $yu_{xx}+u_{yy}=0$.
\end{enumerate}

\item\label{ex:first-order-characteristic-solution} Consider the first-order equation $yu_x+(x^2+1)u_y=0$.
\begin{enumerate}[label=(\alph*)]
\item Determine the characteristics.
\item Show that the most general solution is any function which is constant along characteristics and use this fact to give a formula for the general solution.
\end{enumerate}

\item\label{ex:stokes-ellipticity} The Stokes system in $\mathbb{R}^3$ consists of the equations
\begin{equation}\label{eq:classification-stokes-system}
\Delta\bm{u}-\nabla p=0,\qquad\operatorname{div}\bm{u}=0.
\end{equation}
Here the unknowns are the vector function $\bm{u}:\mathbb{R}^3\to\mathbb{R}^3$ and the scalar function $p:\mathbb{R}^3\to\mathbb{R}$. Show that this system is elliptic.

\item\label{ex:euler-system-type} The Euler equations in $\mathbb{R}^3$ are
\begin{equation}\label{eq:classification-euler-system}
(\bm{u}\cdot\nabla)\bm{u}-\nabla p=0,\qquad\operatorname{div}\bm{u}=0.
\end{equation}
Here the unknowns are the vector function $\bm{u}:\mathbb{R}^3\to\mathbb{R}^3$ and the scalar function $p:\mathbb{R}^3\to\mathbb{R}$. Show that this system is neither elliptic nor hyperbolic.

\item\label{ex:hyperbolicity-two-matrices} Consider the system
\begin{equation}\label{eq:hyperbolic-two-spatial-directions-system}
\bm{u}_t+\bm{A}\bm{u}_x+\bm{B}\bm{u}_y=0.
\end{equation}
What condition must $\bm{A}$ and $\bm{B}$ satisfy for the system to be hyperbolic in the $t$ direction? The condition which you will find is, in general, difficult to verify. Can you give a simple special case? Consider now the special case where $\bm{A}$ and $\bm{B}$ are diagonal. Under what conditions is the system strictly hyperbolic?

\end{enumerate}
